
$b\in K^{\bullet}(\mathbb P(\underline{E}))$, on a
\[
i_*(b\cdot h^*(a))=i_*(b)\cdot g^*(a),
\]
où $g$ et $h$
sont les morphismes structuraux de $\mathbb P(\underline{F})$ et $\mathbb P(\underline{E})$. Comme $h=g\circ i$, on a:
\[
i_*(b\cdot h^*(a))=i_*(b\cdot i^*g^*(a))=g^*(a)\cdot i_*(b)
\]
d'après IV 2.12.4.

D'après (3.4.1), on a donc :
\[
K_\bullet(P)\simeq K_\bullet(S)\otimes_{K^{\bullet}(S)}
\bigl(K^{\bullet}(\mathbb P(\underline{F}))/\operatorname{Im}K^{\bullet}(\mathbb P(\underline{E}))\bigr).
\]
Soient $\xi_F$ et $\xi_E$ les classes des faisceaux inversibles canoniques sur $\mathbb P(\underline{F})$ et $\mathbb P(\underline{E})$. Alors, d'après VI 1.1, $K^{\bullet}(\mathbb P(\underline{F}))$ (resp. $K^{\bullet}(\mathbb P(\underline{E}))$) admet pour base sur $K^{\bullet}(S)$ les éléments $\xi_F^k$ pour $k=0,\ldots,-n$ (resp. $\xi_E^k$ pour $k=0,\ldots,-n+1$). La formule de projection donne :
\[
i_*(\xi_E^k)=i_*(i^*(\xi_F^k))=\xi_F^k i_*(1).
\]
Si $\underline{I}$ est l'idéal de $\mathbb P(\underline{E})$ dans $\mathbb P(\underline{F})$, $\underline{I}$ est un idéal inversible, et on a $i_*(1)=1-I$. Par ailleurs, $\underline{I}$ est défini par l'idéal $J$ de $\mathbb S(\underline{F})$ noyau de $\mathbb S(\underline{F})\to \mathbb S(\underline{E})$, qui est engendré par un sous-Module de $\underline{F}$ isomorphe à $\underline{O}_S$; par suite, $\underline{I}\simeq \underline{O}_{\mathbb P(\underline{F})}(-1)$. L'image de $K^{\bullet}(\mathbb P(\underline{E}))$ est donc le sous-$K^{\bullet}(S)$-module engendré par les $\xi_F^k-\xi_F^{k-1}$ pour $k=0,\ldots,-n+1$. L'homomorphisme $K^{\bullet}(S)\longrightarrow K^{\bullet}(\mathbb P(\underline{F}))/\operatorname{Im}(K^{\bullet}(\mathbb P(\underline{E})))$ est donc un isomorphisme, et par suite, il en est de même pour l'homomorphisme $K_\bullet(S)\to K_\bullet(P)$.

\subsection*{4. Fibrés principaux sous les groupes $\mathbf{Gl}(n)_S$}

\textbf{Proposition 4.1.} \uline{Soient} $S$ \uline{un schéma noethérien}, $G$ \uline{le} $S$-\uline{groupe} $\mathbf{Gl}(n)_S$, $P$ \uline{un torseur sous} $G$, \uline{et} $\underline{E}$ \uline{le} $\underline{O}_S$\uline{-Module localement libre de rang} $n$ \uline{correspondant à} $P$. \uline{On note} $E$ \uline{la classe de} $\underline{E}$ \uline{dans} $K^{\bullet}(S)$, \uline{et} $J$ \uline{le} $\lambda$-\uline{idéal engendré}
